\textbf{Время:} резерв

\textbf{Тезис:} архивный расчёт 2025 года показывает корреляционный избыток в трёх уже введённых направлениях out, side и long.

\textbf{На экране:} сначала шапка «SMASH · расчёт 2025 · Au+Au, 3 ГэВ» и общие отборы внизу; затем три проекции слева направо — out, side, long, в каждой горизонтальная ось и область около q = 0.

\textbf{Текст:} Здесь результаты моей работы 2025 года. В шапке указан расчёт SMASH для столкновений золота при трёх ГэВ. Нижняя подпись задаёт одну выбранную группу: центральность от нуля до десяти процентов и k-T от 0,15 до 0,25 ГэВ на c.

Слева направо показаны out, side и long. По горизонтали в каждой панели — соответствующая компонента разности импульсов, по вертикали — корреляционная функция. Вдали от нуля она близка к единице; около нуля виден избыток парного сигнала. Красные точки и линии — исходный результат и исходная аппроксимация из работы.

Мы узнаём ту же трёхмерную структуру, которую рассматривали для STAR. Каждая картинка показывает одну её проекцию. Ширины в разных направлениях описываются разными радиусами. Сейчас используем архивные графики как пример выполненного анализа.

\textbf{Переход:} Начнём с радиусов двух поперечных направлений — out и side.

\textbf{Справка:} Балинов (2025), исходный слайд 15. Показаны неизменённые растровые панели с исходными осями, точками и линиями; технические ROOT-имена обрезаны. Точное окно усреднения по двум остальным компонентам q и отдельный заряд не указаны; окно ±0.05 ГэВ/c из STAR этим проекциям не приписывается. Высота проекции не обязана равняться двум, поскольку остальные компоненты при усреднении могут быть ненулевыми. Код аппроксимации, её настройки и физическая корректность численного масштаба архивных радиусов независимо не подтверждены. Графики не подменяются результатами или кодом 2026 года.
